% Module 3: Automating Research Tasks
% Beamer Presentation — Claude Code for Academic Researchers
\documentclass[aspectratio=169, 11pt]{beamer}

% ── Theme and Colors ──────────────────────────────────────────────
\usetheme{Madrid}
\usecolortheme{default}

% Custom colors
\definecolor{anthropic}{RGB}{50, 50, 80}
\definecolor{accent}{RGB}{180, 90, 40}
\definecolor{lightgray}{RGB}{245, 245, 245}

\setbeamercolor{structure}{fg=anthropic}
\setbeamercolor{frametitle}{bg=anthropic, fg=white}
\setbeamercolor{title}{bg=anthropic, fg=white}
\setbeamercolor{block title}{bg=anthropic, fg=white}
\setbeamercolor{block body}{bg=lightgray}
\setbeamercolor{alerted text}{fg=accent}

\setbeamerfont{title}{size=\Large, series=\bfseries}
\setbeamerfont{frametitle}{size=\large, series=\bfseries}

% ── Packages ──────────────────────────────────────────────────────
\usepackage{graphicx}
\usepackage{booktabs}
\usepackage{listings}
\usepackage{fontenc}
\usepackage{hyperref}
\usepackage{tikz}

\lstset{
  basicstyle=\ttfamily\small,
  backgroundcolor=\color{lightgray},
  frame=single,
  framerule=0pt,
  breaklines=true,
  columns=fullflexible,
}

% ── Metadata ──────────────────────────────────────────────────────
\title{Module 3: Automating Research Tasks}
\subtitle{Claude Code for Academic Researchers}
\author{} % Instructor name
\date{}
\institute{}

% ══════════════════════════════════════════════════════════════════
\begin{document}

% ── Title ─────────────────────────────────────────────────────────
\begin{frame}[plain]
  \titlepage
\end{frame}

% ══════════════════════════════════════════════════════════════════
%  STRATEGIES FOR WORKING WITH CLAUDE CODE
% ══════════════════════════════════════════════════════════════════

\begin{frame}{Strategies for Working Within the Context Window}
  Recall from Module 2: Claude Code can only ``see'' a limited amount of text at once.\\
  Errors compound when it loses track of earlier decisions.

  \vspace{0.5em}
  \textbf{Five strategies to work within this constraint:}
  \begin{enumerate}
    \item Keep tasks focused
    \item Maintain project documentation
    \item Summarize work between conversations
    \item Validate iteratively
    \item Write token-aware prompts
  \end{enumerate}
\end{frame}

% ── Context Compacting ──────────────────────────────────────────
\begin{frame}{What Happens When the Context Window Fills Up}
  As a conversation grows, Claude Code eventually runs out of room.\\
  When this happens, it \textbf{auto-compacts}: it summarizes older parts of the conversation to free up space.

  \vspace{0.5em}
  \textbf{What auto-compacting does:}
  \begin{itemize}
    \item Replaces earlier messages with a compressed summary
    \item Keeps recent messages and tool results intact
    \item Happens automatically. You will see a notice in the chat
  \end{itemize}

  \vspace{0.5em}
  \textbf{The risk:} Details from early in the conversation (variable definitions, decisions you discussed, constraints you mentioned) may be lost or simplified in the summary.

  \vspace{0.5em}
  \textbf{You can also compact manually:}\\
  Type \texttt{/compact} to force a summary and free up context space.\\
  Useful when you notice Claude Code forgetting earlier decisions or when you want to pivot to a new subtask within the same conversation.

  \vspace{0.5em}
  \alert{This is why CLAUDE.md matters:} it survives compacting because it is reloaded from disk, not from conversation history.
\end{frame}

% ── Strategy 1 ──────────────────────────────────────────────────
\begin{frame}{Strategy 1: Keep Tasks Focused}
  \centering
  {\large \textbf{One topic, one conversation.}}

  \vspace{1em}
  \begin{columns}[T]
    \begin{column}{0.48\textwidth}
      \textbf{\color{red!70!black} Avoid:}
      \begin{itemize}
        \item ``Clean the data, build the panel, run the regressions, and make tables''
        \item Long conversations that drift across topics
        \item Asking Claude Code to hold your entire project in its head
      \end{itemize}
    \end{column}
    \begin{column}{0.48\textwidth}
      \textbf{\color{green!50!black} Instead:}
      \begin{itemize}
        \item ``Clean and reshape the raw census file''
        \item Start a new conversation for each distinct task
        \item Let each conversation do one thing well
      \end{itemize}
    \end{column}
  \end{columns}

  \vspace{0.8em}
  \begin{alertblock}{Why}
    A focused task fits comfortably in the context window. A sprawling task forces Claude Code to forget earlier steps, and that is where silent errors enter.
  \end{alertblock}
\end{frame}

% ── Strategy 2 ──────────────────────────────────────────────────
\begin{frame}{Strategy 2: Maintain Project Documentation}
  Claude Code reads certain files \textbf{automatically} every time it starts.

  \vspace{0.5em}
  \textbf{Keep these up to date:}
  \begin{itemize}
    \item \texttt{CLAUDE.md}: conventions, constraints, variable definitions
    \item \texttt{README.md}: project structure, what each directory contains
    \item Directory-level \texttt{CLAUDE.md} files for subfolders with specific rules
  \end{itemize}

  \vspace{0.5em}
  \textbf{Tell Claude Code what to \textit{ignore}:}
  \begin{itemize}
    \item \texttt{.claudeignore} works like \texttt{.gitignore}
    \item Exclude large data files, build artifacts, logs
    \item Prevents Claude Code from wasting context on irrelevant files
  \end{itemize}

  \vspace{0.5em}
  \begin{block}{The Payoff}
    Good documentation means Claude Code starts every conversation already knowing your conventions, without you spending tokens re-explaining them.
  \end{block}
\end{frame}

% ── Strategy 3 ──────────────────────────────────────────────────
\begin{frame}[fragile]{Strategy 3: Summarize Work Between Conversations}
  When you finish a conversation, the context is gone. The next conversation starts fresh.

  \vspace{0.8em}
  \textbf{Before ending a conversation:}
  \begin{enumerate}
    \item Ask Claude Code to summarize what it did and what remains
    \item Save that summary (in a file, a commit message, or your notes)
    \item Paste the summary into the \textit{next} conversation as context
  \end{enumerate}

  \vspace{0.8em}
  \textbf{Example prompt:}
  \begin{lstlisting}
Summarize what we did in this session: what files
were changed, what decisions were made, and what
still needs to be done.
  \end{lstlisting}

  \vspace{0.5em}
  \alert{Think of it like a handoff memo between research assistants.}\\
  The next ``assistant'' knows nothing unless you brief it.
\end{frame}

% ── Strategy 4 ──────────────────────────────────────────────────
\begin{frame}{Strategy 4: Validate Iteratively}
  \begin{alertblock}{The Compounding Error Problem}
    If step 1 is wrong and you do not catch it, steps 2--5 build on that mistake. By the time you check the final output, the error is buried.
  \end{alertblock}

  \vspace{0.5em}
  \textbf{Check each step before moving to the next:}
  \begin{enumerate}
    \item Clean the data $\rightarrow$ \textit{verify row counts and distributions}
    \item Construct variables $\rightarrow$ \textit{spot-check values against known cases}
    \item Run the merge $\rightarrow$ \textit{confirm pre- and post-merge row counts}
  \end{enumerate}

  \vspace{0.5em}
  Slower per step, but \textbf{much faster overall} because you never have to unwind a chain of errors you didn't catch.
\end{frame}

% ── Strategy 5 ──────────────────────────────────────────────────
\begin{frame}{Strategy 5: Token-Aware Prompting}
  Every word in your prompt costs tokens. Make them count.

  \vspace{0.5em}
  \textbf{A good prompt tells Claude Code three things:}
  \begin{enumerate}
    \item \textbf{What files matter:} ``Read \texttt{data/raw/census.csv} and \texttt{scripts/clean.py}''
    \item \textbf{What the goal is:} ``Create a county-year panel with one row per county per decade''
    \item \textbf{What constraints apply:} ``Keep all 3,143 counties; do not drop rows with missing population''
  \end{enumerate}

  \vspace{0.5em}
  \begin{columns}[T]
    \begin{column}{0.48\textwidth}
      \textbf{\color{red!70!black} Vague:}\\
      {\small ``Clean up the data and make it ready for analysis''}
    \end{column}
    \begin{column}{0.48\textwidth}
      \textbf{\color{green!50!black} Specific:}\\
      {\small ``Read \texttt{census.csv}, drop rows where \texttt{year < 1950}, and save as \texttt{panel\_clean.csv}''}
    \end{column}
  \end{columns}

  \vspace{0.8em}
  \begin{block}{The Rule}
    Be specific about \textit{inputs}, \textit{outputs}, and \textit{constraints}. Be brief about everything else. Context is finite; spend it on what matters.
  \end{block}
\end{frame}

% ══════════════════════════════════════════════════════════════════
%  VALIDATION TECHNIQUES
% ══════════════════════════════════════════════════════════════════

% ── Validation 1 ────────────────────────────────────────────────
\begin{frame}[fragile]{Validation: Check Known Properties}
  Before you look at results, ask: \textit{what do I already know must be true?}

  \vspace{0.5em}
  \textbf{Examples:}
  \begin{itemize}
    \item A 1:1 merge should not change the number of rows
    \item A county-year panel with 50 states $\times$ 10 years = 500 rows
    \item A share variable must be between 0 and 1
    \item No duplicate county-year pairs after reshaping
  \end{itemize}

  \vspace{0.5em}
  \textbf{Ask Claude Code to assert these explicitly:}
  \begin{lstlisting}
After the merge, assert that the output has the
same number of rows as the left table and that
there are no duplicate fips-year pairs.
  \end{lstlisting}

  If the assertion fails, you catch the error immediately, not three steps later.
\end{frame}

% ── Validation 2 ────────────────────────────────────────────────
\begin{frame}{Validation: Write Down Expectations First}
  \begin{block}{The Pre-Registration Analogy}
    Just as you pre-register hypotheses before seeing results, write down what a correct transformation should produce \textit{before} you run it.
  \end{block}

  \vspace{0.5em}
  \textbf{Before a merge or transformation, state:}
  \begin{itemize}
    \item Expected output dimensions (rows $\times$ columns)
    \item Which columns should appear and their types
    \item How missing values should be handled
    \item A specific row you can trace through by hand
  \end{itemize}

  \vspace{0.5em}
  \textbf{Then compare the actual result to your expectations.}\\
  Discrepancies are either a bug or a misunderstanding, both worth catching early.
\end{frame}

% ── Validation 3 ────────────────────────────────────────────────
\begin{frame}[fragile]{Validation: Unit Tests Against Known Values}
  Pick a handful of observations and compute the correct answer by hand.\\
  Then verify the code reproduces them.

  \vspace{0.5em}
  \textbf{Example prompt:}
  \begin{lstlisting}
Write a test that verifies:
- Cook County (fips=17031) in 2010 has
  population = 5,194,675
- The migration_share for fips=06037
  in 2000 is between 0.28 and 0.32
  \end{lstlisting}

  \vspace{0.5em}
  \textbf{Why this works:}
  \begin{itemize}
    \item Forces you to engage with the data at the observation level
    \item Catches silent recoding, wrong joins, and off-by-one errors
    \item Tests survive across refactors. Rerun them after every change
  \end{itemize}
\end{frame}

% ── Validation 4 ────────────────────────────────────────────────
\begin{frame}{Validation: Replicate Before You Extend}
  \textbf{Building on someone else's work?}\\
  Have Claude Code reproduce their results first.

  \vspace{0.5em}
  \begin{enumerate}
    \item Give Claude Code the original paper's code or description
    \item Ask it to replicate a key table or figure
    \item Compare output to the published result
    \item Only after replication succeeds, begin your extension
  \end{enumerate}

  \vspace{0.5em}
  \begin{alertblock}{Why}
    If you cannot reproduce the baseline, anything you build on top of it is uninterpretable. Replication is not a detour; it is the foundation.
  \end{alertblock}
\end{frame}

% ── Validation 5 ────────────────────────────────────────────────
\begin{frame}{Validation: Adversarial Agents}
  Use a \textit{second} Claude Code conversation as a critic of the first.

  \vspace{0.5em}
  \textbf{How:}
  \begin{enumerate}
    \item Agent 1 writes the code
    \item Open a new conversation (Agent 2) and ask it to review:\\
      {\small ``Read \texttt{scripts/build\_panel.py} and identify any bugs, silent data loss, or decisions that could change the results.''}
    \item Agent 2 has fresh context and no attachment to the code
  \end{enumerate}

  \vspace{0.5em}
  \textbf{Why this is powerful:}
  \begin{itemize}
    \item The writing agent has blind spots because it wrote the code and ``believes'' it is correct
    \item A fresh agent reads the code like a skeptical reviewer
    \item Catches issues you would miss because you watched it being built
  \end{itemize}

  \vspace{0.5em}
  \alert{Treat your agents like co-authors who check each other's work, not a single assistant you trust unconditionally.}
\end{frame}

% TODO: Add exercise slides, debrief, and summary for Module 3

\end{document}
